\section{Method: A Design Space of Physical Injection}
\label{sec:method}

We answer the question of \S\ref{sec:intro} by construction: design the ways physics can be injected into a shared latent, state what each design is supposed to achieve, and lay out the experiments that test the mechanism behind it.

\subsection{Injection Design Space: Five Families, Ten Variants}
\label{sec:designspace}

We span the design space along three orthogonal axes---\emph{hard vs.\ soft} constraints on the physical quantity, pinning \emph{what the motion state is} vs.\ \emph{how it evolves}, and \emph{labeled vs.\ label-free} supervision---so that any ``try a different encoding'' objection lands on an axis we have measured. Table~\ref{tab:mechanisms} lists the ten variants with the design rationale of each.

\begin{table}[tb]
\centering
\small
\setlength{\tabcolsep}{3pt}
\begin{tabular}{@{}p{2.35cm}p{5.4cm}@{}}
\toprule
Variant & What it does, and what it tests \\
\midrule
\multicolumn{2}{@{}l}{\emph{Family 1: fixed slot --- hard constraint, pins the state, labeled}} \\
\ \ slot (structpos) & pin dims $[0{:}2]$ to the true position (the object's $x,y$ coordinates): the most direct state injection; tests whether a hard constraint helps \\
\ \ \ $+$reweight & raise the slot's share of the prediction loss (weight $30$): if failure is due to the slot's ${\sim}1\%$ gradient share, this should rescue it \\
\ \ \ $+$velocity & encode velocity too: tests whether it matters \emph{which} quantity is injected, not just that something is \\
\multicolumn{2}{@{}l}{\emph{Family 2: deep-supervision probe --- soft constraint, pins the state, labeled}} \\
\ \ probe & linear head reads position from the latent: replicates the deep-supervision recipe \citep{zahorodnii2025deepsup} in a high-dimensional visual latent \\
\ \ \ $+$slot & soft $+$ hard combination: tests complementarity \\
\multicolumn{2}{@{}l}{\emph{Family 3: kinematic head --- hard constraint, pins the evolution, labeled}} \\
\ \ free MLP & attach $z{+}v{+}a$ update with learned $a$: upgrade from pinning state to pinning \emph{evolution} \\
\ \ strict $a{=}g$ & exact gravitational form, one learnable constant: closes the ``your physics was wrong'' objection \\
\multicolumn{2}{@{}l}{\emph{Family 4: consistency --- soft constraint, pins the evolution, labeled}} \\
\ \ consistency & finite-difference velocity of the rolled-out positions must match ground truth; assumes no form for $a$, designed for collision impulses \\
\multicolumn{2}{@{}l}{\emph{Family 5: label-free prior --- soft constraint, pins the evolution, label-free}} \\
\ \ label-free & slot must evolve as smooth second-order dynamics, no ground truth: the only injection available for raw video \\
\ \ grounded & same structure $+$ true labels: isolates the label variable \\
\bottomrule
\end{tabular}
\caption{Ten injection variants in five families, spanning three axes (read off each family's header). \textbf{Constraint}---\emph{hard}: a latent dimension is forced to equal the physical quantity; \emph{soft}: an auxiliary loss only encourages it. \textbf{Target}---the \emph{state} (what the motion is: position, velocity) vs.\ its \emph{evolution} (how it changes: a dynamics rule). \textbf{Supervision}---\emph{labeled} uses ground-truth physics, \emph{label-free} does not. The three axes span the ways to inject physics into a shared latent, so any alternative encoding falls on an axis we test. All variants use the strongest training protocol (\S\ref{sec:protocol}).}
\label{tab:mechanisms}
\end{table}

Orthogonally, we cross injection with the \emph{training regime}: fine-tuning a pretrained encoder vs.\ from-scratch co-training; with injection on/off this gives a $2{\times}2$ per domain. This closes the remaining escape hatch---``structure must be baked in during pretraining, as extrinsic architectures do.''

\subsection{Mechanistic Hypothesis and Its Tests}
\label{sec:mechtests}

\emph{Hypothesis (the load-bearing problem).} The physical slot occupies $2$ of $192$ dimensions and receives ${\sim}1\%$ of the prediction gradient, while the black-box channel \emph{redundantly} encodes the same position---so prediction can route around any injected slot. Because the latent already carries the injected state, injection adds no new signal---it only consumes capacity and gradient share. We test this account in three ways, each admitting a specific outcome that would falsify it:

\begin{enumerate}
    \item \textbf{Is the physical constraint injected?} Split the latent into the $2$ slot dimensions and the $190$ black-box dimensions and decode position from each separately, with a random-2-dimension control. If the black box alone cannot decode position, the redundancy claim is false.
    \item \textbf{Is the constraint acting at all?} Two measurements: how large the weighted physical auxiliary loss is relative to the prediction loss (a proxy for gradient magnitude), and whether the encoder's effective dimensionality (participation ratio) changes. If neither moves, the constraint never took effect and the failure is a bug. If the auxiliary term dominates and dimensionality collapses, the constraint is acting---and competing with prediction.
    \item \textbf{Does re-weighting the physical loss help?} Sweep the physical slot's prediction-loss weight from $1$ to $300$. If harm does not respond systematically to the weight, the gradient-share mechanism is wrong; if it responds but weighting to saturation still cannot rescue every domain, demonstrating that insufficient physical loss is not the root cause.
\end{enumerate}

\subsection{The Training-Protocol Control: Free Rollout}
\label{sec:protocol}

Free rollout serves two methodological roles; neither is a novelty claim (it is scheduled sampling, \citealp{bengio2015scheduled}).

\emph{Controlled factorial variable.} Physical structure has never been compared against the training protocol on the same baseline; prior studies treat the protocol as background. We set the protocol first---replacing teacher forcing (one-step, ground-truth context) with free rollout ($H{=}8$ autoregressive steps, supervising the whole rollout)---and then ask whether structure adds incremental value. Teacher forcing hides compounding error at training time (\emph{exposure bias}): the model never sees its own mistakes, then must consume them at deployment. Free rollout exposes them during training, injecting no physics whatsoever. We read the resulting gain as more than robustness to compounding error: supervising an $H$-step autoregressive rollout turns training itself into \emph{long-horizon dynamics simulation}. Under teacher forcing the predictor is graded only on reproducing local one-step transitions; under free rollout it is graded on \emph{evolving} the state over the full horizon, so the gradient rewards a self-consistent multi-step dynamics rather than a one-step map---precisely the capability the latent is missing (\S\ref{sec:intro}). On this reading, free rollout is an evolution-targeted training signal that requires no physics labels: it trains \emph{how the state evolves} without touching \emph{what the state is}. A second protocol variable, the rollout horizon $H$, must \emph{match dynamics complexity}: domains with late causal events (collision; realistic-simulation scenes) need horizons long enough to contain them.

\emph{Positive control.} A negative result invites two deflations---``your implementation is broken'' or ``your metrics cannot detect improvement.'' A protocol change that produces large gains on identical code and identical metrics rules out both at once, pinning the failure on the injected structure itself.

\subsection{External and Cross-Backbone Controls}
\label{sec:external}

\emph{Extrinsic port.} We faithfully port the physically interpretable world model of \citet{mao2024piwm}---a three-stage pipeline (VAE encoder, $128$-D; a supervised extractor to a low-dimensional physical state; known-equation dynamics with learnable constants)---onto the same PhyWorld domains and OOD protocol. The port tests our mechanism's architectural prediction from the outside: extrinsic designs make the physical state the sole pathway prediction must route through, closing the bypass by construction.

\emph{Cross-backbone control.} To test that the conclusions are not artifacts of one encoder, we replicate the core protocol on a second JEPA-family instantiation in the style of DINO-WM \citep{zhou2025dinowm}: a \emph{frozen} DINOv2-small encoder \citep{oquab2024dinov2}---generic self-supervised features that have never seen PhyWorld---with the trainable projector serving as adapter and the identical predictor stack, losses, seeds, and evaluation. The variant differs from LeWM in backbone provenance, dimensionality, and, crucially, in whether the injected losses can reshape the encoder at all.
